% Chapter 4 — Integrating LichtFeld Studio v0.5.2
% Video-to-Gaussian: An Agentic Pipeline for Volumetric Cultural Heritage Preservation
% Dr John O'Hare, DreamLab AI Consulting Ltd
% v2.0 — June 2026

\chapter{Integrating LichtFeld Studio v0.5.2}
\label{ch:lichtfeld}

The v1.0 pilot proved that \acrshort{3dgs}~\parencite{Kerbl2023} was a viable medium for heritage preservation.
Its bespoke trainer, however, was a maintenance liability: every advance in densification strategy, rasteriser
performance, or export format required hand-porting rather than a simple upstream pull.
This chapter documents how the v2.0 system was re-platformed onto
LichtFeld~Studio~v0.5.2~\parencite{LichtFeldStudio2026} (git tag \texttt{d9e76009}, dated 2026-04-21)
and describes the specific capabilities that upgrade delivered to the heritage capture workflow.

% ─────────────────────────────────────────────────────────────────────────────
\section{Rationale for Re-Platforming}
\label{sec:lf-rationale}
% ─────────────────────────────────────────────────────────────────────────────

The v1.0 bespoke trainer grew organically from an early open-source \acrshort{3dgs} reference implementation.
Over approximately eighteen months it accumulated local patches for multi-resolution densification, a
custom CUDA rasteriser backport, and ad-hoc export routines for \texttt{.ply} and \texttt{.splat} containers.
Each of those additions was written against internal data structures that diverged progressively from the
community codebase, making it increasingly difficult to absorb improvements from the research stream
without substantial merge effort.

LichtFeld~Studio emerged as the community's de-facto maintained distribution of \acrshort{3dgs} tooling.
By v0.5.2 it combined a rigorously benchmarked CUDA rasteriser, a non-destructive editing suite, native
\acrshort{usd} I/O, an in-application quality-evaluation mode, and a programmatic \acrshort{mcp} control
surface — each component maintained by a community with a faster release cadence than a single-institution
project can sustain.

The re-platforming decision was therefore primarily an exercise in reducing bespoke code surface.
The \texttt{gaussian-toolkit} fork tracks 215 upstream commits to the v0.5.2 tag and adds only the
integration layer described in this chapter: scene-graph mapping, orchestrator bindings, and the
quality-gating logic that consumes evaluation metrics from the engine.
All reconstruction logic, the CUDA rasteriser, densification strategies (\acrshort{mcmc},
\acrshort{mrnf}), and mesh-extraction backend stubs reside in the upstream codebase.
This arrangement allows the toolkit to absorb research velocity from the LichtFeld community
while confining heritage-specific concerns to a thin, reviewable integration layer.

% ─────────────────────────────────────────────────────────────────────────────
\section{Scene-Graph Harmonisation}
\label{sec:lf-scenegraph}
% ─────────────────────────────────────────────────────────────────────────────

Upstream commit \texttt{34e2ba85} introduced what the LichtFeld project calls
\emph{scene-graph harmonisation}: a unified internal node hierarchy that supersedes the earlier
flat list of Gaussian primitive arrays.
Prior to that commit, the trainer maintained Gaussians as unstructured parameter tensors with no
persistent identity across training iterations; exporting a scene meant serialising the full tensor
state into a monolithic \texttt{.ply} file.
After harmonisation, every Gaussian is assigned to a named node in an acyclic hierarchy, and the
hierarchy itself is preserved across checkpoint saves and resumes.

The toolkit's integration layer maps this harmonised node hierarchy onto \acrshort{usd}
prims~\parencite{Pixar2019} at export time using a deterministic naming convention:
each internal node identifier becomes a USD \texttt{Xform} prim path, and the Gaussian
payload attached to that node becomes a nested \texttt{PointInstancer} or
\texttt{GaussianSplat} schema prim depending on downstream consumer preference.
The result is a clean, one-to-one correspondence between the trainer's internal scene structure
and the exported USD scene graph.

This correspondence is the key enabler of the hierarchical export described in
Chapter~\ref{ch:scenegraph}.
Because the mapping is defined by node names rather than tensor indices, it is stable across
re-training runs and amenable to manual curation: a conservator may rename or reparent nodes in the
USD file and re-import them into the trainer without corruption.
The scene graph therefore functions as the shared address space that binds the training engine,
the editing suite, and the downstream assembly pipeline into a single coherent workflow.

% ─────────────────────────────────────────────────────────────────────────────
\section{Native USD I/O}
\label{sec:lf-usd}
% ─────────────────────────────────────────────────────────────────────────────

LichtFeld~Studio~v0.5.2 introduced a first-class \acrshort{usd} layer built on the
OpenUSD C++ libraries~\parencite{Pixar2019}, exposing bidirectional scene
read and write from within the training process itself.
Earlier workflows required a post-training conversion step: the trainer emitted \texttt{.ply},
a separate tool converted to USD, and any editing performed in a USD-aware host had to be
re-imported through a lossy round-trip.

Under the native I/O model the trainer reads initial point-cloud priors directly from a USD layer,
preserves camera rig geometry as USD camera prims, and checkpoints the Gaussian field into a
\texttt{.usda}/\texttt{.usdz} stage at each evaluation interval.
The editing suite — responsible for splat culling, object segmentation, and manual correction —
opens the same USD stage, performs non-destructive overrides via USD's \emph{composition arcs},
and saves back without touching the base layer written by the trainer.
The assembly pipeline described in Chapter~\ref{ch:scenegraph} then aggregates per-object
USD stages via reference and payload arcs, producing the final hierarchical scene without any
format conversion at any stage of the process.

This end-to-end USD round-trip is significant for heritage work specifically because it eliminates
the provenance ambiguity introduced by intermediate formats.
Every transformation applied to the scene is recorded as a USD opinion on a named layer with a
defined composition strength, making the full edit history auditable.

% ─────────────────────────────────────────────────────────────────────────────
\section{GUI Evaluation Mode}
\label{sec:lf-eval}
% ─────────────────────────────────────────────────────────────────────────────

LichtFeld~Studio v0.5.2 includes an in-application held-out-view evaluation mode that reports
image-quality metrics continuously during and after training.
The evaluation set is defined by withholding a configurable fraction of the reconstructed camera
poses before training begins; the trainer renders from those withheld viewpoints at each evaluation
step and computes three metrics against the corresponding source frames:

\begin{itemize}
  \item \textbf{\acrshort{psnr}} — the classical pixel-fidelity measure, expressed in decibels;
        higher is better.
  \item \textbf{\acrshort{ssim}}~\parencite{Wang2004SSIM} — a perceptual measure of luminance,
        contrast, and structural correlation; range $[0,1]$, higher is better.
  \item \textbf{\acrshort{lpips}}~\parencite{Zhang2018LPIPS} — a deep-feature perceptual
        distance; lower is better.
\end{itemize}

The GUI presents these as live plots during training, enabling an operator to observe convergence
and abort an under-performing run early.
After training, the same mode produces a per-scene evaluation report that the toolkit consumes
for automated quality gating: a training run is promoted to the export stage only when
\acrshort{psnr} exceeds 28\,dB, \acrshort{ssim} exceeds 0.90, and \acrshort{lpips} falls below
0.15.

For the reference gallery scene described in Chapter~\ref{ch:reconstruction}, GUI Evaluation Mode
reported \textbf{PSNR 29.6\,dB}, \textbf{SSIM 0.921}, and \textbf{LPIPS 0.108}, all comfortably
within the quality gate.
Figure~\ref{fig:lf-eval} shows a representative held-out-view render produced by the evaluation
pass.
Full per-scene metric tables and convergence plots are presented in Chapter~\ref{ch:reconstruction}.

\begin{figure}[htbp]
  \centering
  \includegraphics[width=0.85\textwidth]{v2/splat_render.png}
  \caption{A rendered view from the trained Gaussian field, evaluated in GUI Eval Mode.
           The held-out viewpoint was not seen during training; the rendered image achieved
           PSNR~29.6\,dB, SSIM~0.921, LPIPS~0.108 against the source frame.}
  \label{fig:lf-eval}
\end{figure}

% ─────────────────────────────────────────────────────────────────────────────
\section{Model Context Protocol Control Surface}
\label{sec:lf-mcp}
% ─────────────────────────────────────────────────────────────────────────────

The defining architectural feature of LichtFeld~Studio~v0.5.2 for this project is its built-in
\acrshort{mcp} server.
At startup the engine binds a JSON-RPC endpoint on port~45677 and advertises a catalogue of
70$+$ callable tools organised into the functional categories summarised in
Table~\ref{tab:mcp-tools}.
The agentic orchestrator integrated into the \texttt{gaussian-toolkit} drives the entire
training, evaluation, and export cycle exclusively through this interface — no shell scraping,
no log parsing, no filesystem polling.

\begin{table}[htbp]
  \centering
  \caption{Representative \acrshort{mcp} tool categories exposed by LichtFeld~Studio on port~45677.}
  \label{tab:mcp-tools}
  \small
  \begin{tabular}{@{}lp{8.8cm}@{}}
    \toprule
    \textbf{Category} & \textbf{Representative tools} \\
    \midrule
    Training control
      & \texttt{training.start}, \texttt{training.stop}, \texttt{training.set\_iterations},
        \texttt{training.set\_densification\_strategy} \\[4pt]
    Checkpoint \& evaluation
      & \texttt{checkpoint.save}, \texttt{checkpoint.restore},
        \texttt{eval.run\_held\_out}, \texttt{eval.get\_metrics} \\[4pt]
    Camera \& pose
      & \texttt{cameras.list}, \texttt{cameras.import\_colmap},
        \texttt{poseopt.enable}, \texttt{poseopt.get\_residuals} \\[4pt]
    Export \& USD
      & \texttt{export.usd}, \texttt{export.ply}, \texttt{export.splat},
        \texttt{usd.set\_layer\_metadata} \\[4pt]
    Scene-graph queries
      & \texttt{scene.list\_nodes}, \texttt{scene.get\_node\_bbox},
        \texttt{scene.reparent}, \texttt{scene.prune\_by\_opacity} \\
    \bottomrule
  \end{tabular}
\end{table}

The integration seam between the toolkit's orchestrator and the LichtFeld engine is illustrated
in Figure~\ref{fig:mcp-seam}.
The orchestrator dispatches tool calls; the \acrshort{mcp} server translates them into engine API
calls; the engine returns structured JSON responses.
This separation of concerns means the orchestrator never depends on the internal implementation
of the trainer — it interacts only with the stable tool catalogue, which the LichtFeld project
treats as a versioned public API.
Tool responses include typed payloads (e.g.\ metric dictionaries, node identifier lists, USD
layer paths) that the orchestrator uses for branching decisions without parsing free-form text.

\begin{figure}[htbp]
  \centering
  \begin{tikzpicture}[
    node distance=1.6cm and 2.2cm,
    box/.style={
      rectangle, rounded corners=4pt,
      draw=#1!70, fill=#1!12,
      text width=3.4cm, minimum height=1.1cm,
      align=center, font=\small\sffamily
    },
    arr/.style={-{Stealth[length=3.5mm]}, thick, MidGrey},
    lbl/.style={font=\footnotesize\sffamily, MidGrey}
  ]
    \node[box=DreamLabAccent] (orch)
      {\textbf{Agentic}\\Orchestrator};
    \node[box=UoSRed, right=of orch] (mcp)
      {\textbf{MCP Server}\\port 45677};
    \node[box=MidGrey, right=of mcp] (eng)
      {\textbf{LichtFeld}\\Engine};

    \draw[arr] ([yshift= 3pt]orch.east) -- ([yshift= 3pt]mcp.west)
      node[midway, above, lbl] {tool call};
    \draw[arr] ([yshift=-3pt]mcp.west) -- ([yshift=-3pt]orch.east)
      node[midway, below, lbl] {JSON response};

    \draw[arr] ([yshift= 3pt]mcp.east) -- ([yshift= 3pt]eng.west)
      node[midway, above, lbl] {API call};
    \draw[arr] ([yshift=-3pt]eng.west) -- ([yshift=-3pt]mcp.east)
      node[midway, below, lbl] {typed payload};
  \end{tikzpicture}
  \caption{Integration seam between the \texttt{gaussian-toolkit} orchestrator and
           LichtFeld~Studio~v0.5.2. All control and data flow passes through the
           \acrshort{mcp} server; neither the orchestrator nor downstream pipeline stages
           depend on engine internals.}
  \label{fig:mcp-seam}
\end{figure}

The practical consequence for the heritage workflow is that the entire processing cycle —
initiating training from a \acrshort{colmap} reconstruction, monitoring convergence, triggering
held-out evaluation, gating on quality metrics, exporting a USD stage, and handing off to the
scene-graph assembly stage — is expressed as a deterministic sequence of \acrshort{mcp} tool
calls in the orchestrator's workflow definition.
Re-running the pipeline from any checkpoint requires only re-issuing the corresponding tool
calls; there is no fragile shell interaction to re-create.

% ─────────────────────────────────────────────────────────────────────────────
\section{Supporting Features}
\label{sec:lf-supporting}
% ─────────────────────────────────────────────────────────────────────────────

Three further features of LichtFeld~Studio~v0.5.2 support the heritage capture workflow
beyond the core training and export path.

\paragraph{Pose optimisation (\texttt{poseopt}).}
Heritage captures are frequently taken with hand-held cameras moving through spaces that combine
flat reflective surfaces, repetitive stonework, and extreme contrast ratios.
These conditions are hostile to \acrshort{sfm} pose estimation: \acrshort{colmap} often produces
trajectories with residual rotational drift across long corridors or drift accumulated over a slow
360° sweep of a single object.
LichtFeld's \texttt{poseopt} module, exposed through the \texttt{poseopt.enable} and
\texttt{poseopt.get\_residuals} \acrshort{mcp} tools, jointly optimises camera extrinsics
alongside Gaussian parameters during training using a photometric reprojection loss.
For the reference gallery scene, enabling \texttt{poseopt} reduced mean angular reprojection
error from 1.7\,mrad to 0.4\,mrad over the first 10\,000 training iterations, recovering fine
structural detail in the rear alcoves that was blurred in the rigid-pose baseline.

\paragraph{Timelapse capture.}
LichtFeld emits a rendered viewpoint snapshot at each evaluation step and accumulates these into
an ordered image sequence addressable via the \texttt{eval.get\_timelapse} \acrshort{mcp} tool.
The toolkit assembles these snapshots into an MP4 timelapse that documents training convergence
from the initial noisy point cloud to the final photorealistic field.
This artefact has proved useful for stakeholder communication — conservators unfamiliar with
Gaussian splatting can observe the model ``coming into focus'' over minutes of screen time —
and for post-hoc diagnosis of runs that converged to a local optimum.

\paragraph{\texttt{gaussian-toolkit} utilities (GUT).}
The integration layer ships a small Python library, \texttt{gut}, containing utilities for
checkpoint provenance stamping, metric logging to the project's observability stack, and
orchestrator-side validation of USD export integrity.
These utilities sit entirely outside the LichtFeld codebase and impose no upstream dependency,
ensuring the separation between the engine and the heritage-specific integration layer is
maintained cleanly.

% ─────────────────────────────────────────────────────────────────────────────
\section{Upstream Provenance and Licence Posture}
\label{sec:lf-provenance}
% ─────────────────────────────────────────────────────────────────────────────

The \texttt{gaussian-toolkit} fork incorporates 215 upstream commits from the LichtFeld~Studio
repository up to the v0.5.2 tag (\texttt{d9e76009}), plus the toolkit's own integration
layer described throughout this chapter.
The integration layer comprises approximately 3\,200 lines of Python, organised as the
orchestrator bindings (\texttt{gut.orchestrator}), the USD mapping module
(\texttt{gut.usd\_bridge}), and the quality-gate logic (\texttt{gut.quality}).

Regarding licence posture: the LichtFeld~Studio engine itself is released under an open-source
licence that permits academic and commercial use without restriction.
Two of the four mesh-extraction backends bundled with the toolkit carry additional licence
constraints.
The \acrshort{come}~\parencite{Radl2025CoMe} backend is distributed under the Inria
non-commercial research licence, restricting deployment to non-commercial institutional use.
The \acrshort{gw} backend~\parencite{Waczynska2024GaussianWrapping} carries a comparable
academic-use restriction.
Both restrictions were assessed in \acrshort{adr}-004 and \acrshort{adr}-005 respectively, which
concluded that University of Salford deployment for heritage preservation constitutes qualifying
non-commercial academic use.
The \acrshort{milo}~\parencite{Guedon2025MILo} and \acrshort{tsdf} backends carry no additional
constraints beyond their upstream licences.
Full licence attribution and the text of each \acrshort{adr} are included as appendices to this
report.
The meshing backends and their respective trade-offs are evaluated empirically in
Chapter~\ref{ch:meshing}.
